\chapter{A theorem you can rotate}
\chaptertag{Necklace colorings and Fermat's little theorem}

Why should a prime $p$ divide $a^p-a$ for every integer $a$? Instead of manipulating the expression, we count objects that naturally arrive in packets of $p$.

\section{Words on a circular loom}

Fix a prime $p$ and $a$ colors. First imagine $p$ labeled positions around a circle. A coloring is a word
\[
(c_0,c_1,\dots,c_{p-1}),
\]
and there are $a^p$ such words. Rotation by one step sends it to
\[
(c_{p-1},c_0,\dots,c_{p-2}).
\]

\begin{center}
\begin{tikzpicture}[scale=.82]
\foreach \shift/\rot/\lab in {0/0/{first},4.4/72/{rotate},8.8/144/{rotate again}}{
  \begin{scope}[xshift=\shift cm,rotate=\rot]
    \draw[navy!25,thick] (0,0) circle (1.25);
    \foreach \ang/\col in {90/teal,18/gold,-54/teal,-126/coral,162/gold}{
      \fill[\col] (\ang:1.25) circle (5.5pt);
    }
  \end{scope}
  \node[muted] at (\shift,-1.75) {\lab};
}
\draw[-{Stealth},teal,thick] (1.5,0)--(2.9,0);
\draw[-{Stealth},teal,thick] (5.9,0)--(7.3,0);
\end{tikzpicture}
\end{center}

Group words together when one is a rotation of another. Such a group is called a rotational \emph{orbit}. How large can an orbit be?

\begin{lemma}[Prime orbit dichotomy]
For $p$ prime, a length-$p$ word has either $1$ or $p$ distinct rotations. It has one exactly when it is constant.
\end{lemma}

\begin{proof}
If rotation by $k$ positions, with $1\le k<p$, fixes the word, then entries $i$ and $i+k$ have the same color (indices modulo $p$). Since $\gcd(k,p)=1$, repeatedly adding $k$ visits every position. Thus all entries have one color. A nonconstant word has no repeated rotation before the $p$th, so its orbit has size $p$.
\end{proof}

There are $a$ constant words, one for each color. The remaining $a^p-a$ words split into orbits of exactly $p$ words.

\begin{theorem}[Fermat's little theorem, necklace form]
If $p$ is prime, then for every integer $a$,
\[
a^p\congr a\pmod p.
\]
If $p\ndivides a$, then $a^{p-1}\congr1\pmod p$.
\end{theorem}

\begin{proof}
For $a\ge0$, the orbit count above shows $p\divides a^p-a$. This congruence depends only on the residue of $a$ modulo $p$, so it holds for every integer $a$. If $p\ndivides a$, then $a$ has an inverse modulo $p$, and we may cancel it from $a^p\congr a$.
\end{proof}

\begin{bigidea}
Divisibility by $p$ appears because symmetry packages nonexceptional objects into groups of $p$. The exceptions are exactly the colorings fixed by rotation. This pattern -- count, isolate fixed points, divide the rest into orbits -- reaches far beyond number theory.
\end{bigidea}

\section{A subtle choice: labeled words or necklaces?}

A physical necklace usually ignores where position $0$ is. Here we first count \emph{labeled} circular words, then use rotation to partition them. Calling each orbit a necklace is convenient, but the proof depends on keeping those two levels distinct.

\begin{warningbox}
When the length is composite, an orbit can have an intermediate size. The word $RBRB$ has two distinct rotations, not $1$ or $4$. Primality is doing real work.
\end{warningbox}

\section{Using the theorem}

To compute $3^{1000}\pmod7$, note that $3^6\congr1$ and $1000=6\cdot166+4$:
\[
3^{1000}\congr(3^6)^{166}3^4\congr81\congr4\pmod7.
\]

Fermat's theorem also explains the first layer of binomial divisibility. For $1\le k\le p-1$,
\[
\binom pk=\frac{p!}{k!(p-k)!}
\]
is divisible by $p$, because the denominator contains no factor $p$. Therefore
\[
(x+y)^p\congr x^p+y^p\pmod p.
\]
We will return to this ``freshman's dream'' after developing binomial coefficients systematically.

\begin{exercisebox}
\begin{enumerate}
  \item List the rotational orbits of binary words of length $3$. Identify the constant words. \exlevel{Fluency}
  \item Compute $11^{2026}\pmod{17}$ using Fermat's theorem. \exlevel{Fluency}
  \item In the orbit lemma, prove carefully that $0,k,2k,\dots,(p-1)k$ are all distinct modulo $p$. \exlevel{Reasoning}
  \item Give an orbit-counting proof that $p\divides a^p-a$ using linear strings and cyclic shift, without drawing a circle. \exlevel{Reasoning}
  \item Find every possible orbit size for binary words of length $6$, and give an example of each. \exlevel{Exploration}
  \item Prove that if $p$ is prime, then $(x+y)^p\congr x^p+y^p\pmod p$. Use induction to deduce $(x_1+\cdots+x_m)^p\congr x_1^p+\cdots+x_m^p$. \exlevel{Reasoning}
  \item If $n$ is composite, must there be an $a$ with $a^n\not\congr a\pmod n$? Test $n=341$ and $561$. This opens the door to pseudoprimes and Carmichael numbers. \exlevel{Challenge}
\end{enumerate}
\end{exercisebox}

